\section{Conclusion}
\label{sec:conclusion}

We presented a controlled empirical study of call-graph-guided LLM fault
localisation on the \sqlite~3.30.1 codebase.
Across nine conditions and two models, graph context provided only marginal
exact-accuracy gains over the no-context baseline.
The primary ceiling is structural: 31\,\% of \qwen{} failure slots arise from
function-pointer gaps that make the target unreachable in both graphs, and
removing these two structurally-blocked questions does not substantially change
the picture for the remaining four.
Where the target is reachable, the dominant failure mode is incorrect function
identification (52\,\% of \qwen{} slots), suggesting that graph traversal alone
is insufficient without better disambiguation among structurally adjacent
functions.
Text slices match the tool condition at ten to fifteen times lower token cost,
making them a competitive practical choice for shallow-target questions.
Caller/callee tools reduce unsupported edge citations by a factor of 2.4
(28\,\% vs.\ 12\,\%) even when they do not change the final answer, indicating a
navigational grounding effect that may matter more at scale.

Future work should address function-pointer recovery to raise the structural
ceiling, expand the benchmark beyond six questions and one codebase, and
investigate whether larger context windows or chain-of-thought prompting change
the tool-vs-text trade-off.
